\SourceLeaf{372}%
\begin{center}
{\large NACHLASS UND BRIEFWECHSEL.\par}
\bigskip
ALLGEMEINSTE REGELN ZUR BESTIMMUNG DER MITTELPUNKTE VON KRÜMMUNGSKREISEN\\
IN JEDEM GEGEBENEN PUNKT EINER GEGEBENEN KURVE.
\end{center}

Die Lage eines beliebigen Punkts in einer gegebenen Ebene sei durch zwei Veränderliche $x,y$ bestimmt (etwa Abstände von gegebenen Punkten oder Geraden, Winkel usw.). Es seien
\[
P=0,\qquad P'=0
\]
die Gleichungen zweier Kurven, wobei $P,P'$ beliebige Funktionen von $x,y$ bezeichnen. Ferner gelte
\[
\begin{aligned}
dP&=p\,dx+q\,dy,&dP'&=p'\,dx+q'\,dy,\\
d^2P&=r\,dx^2+2s\,dx\,dy+t\,dy^2,&d^2P'&=r'\,dx^2+2s'\,dx\,dy+t'\,dy^2.
\end{aligned}
\]
Dann haben die Kurven den durch
\[
x=\xi,\qquad y=\eta
\]
bestimmten Punkt gemeinsam, wenn für diese Werte gilt
\[
P=P';
\]
und in diesem Punkt dieselbe Tangente, wenn außerdem gilt
\[
pq'=p'q;
\]
